\chapter{方法实现}\label{chap:float}

\section{统一状态表示与合法性判定}

六种方法共享同一种状态定义：长度为 $k$ 的元组
\begin{equation}
  s=(q_0,\ldots,q_{k-1})
\end{equation}
表示前 $k$ 列已完成放置。尝试在下一列放入候选行 $r$ 时，判定函数为
\begin{equation}
  \text{safe}(s,r)=\bigwedge_{c=0}^{k-1}\left[(q_c\neq r)\land(|q_c-r|\neq |c-k|)\right].
\end{equation}
若条件成立，状态扩展为 $s'=s\oplus r$。

\section{M1：命令式穷举 + 逻辑筛选}

M1 先由 \texttt{itertools.permutations} 生成全部 $8!$ 个列-行映射，再交给 Kanren 的 \texttt{membero} 与对角线关系做二次筛选。实现直观，但候选空间固定，无法提前剪枝。

\section{M2：纯逻辑求解（\texttt{permuteq}）}

M2 用关系 \texttt{permuteq(rows, qs)} 描述“\texttt{qs} 是 \texttt{rows} 的一个排列”，随后叠加对角线约束：
\begin{equation}
  \text{solution}(\mathbf{q})\Leftarrow \texttt{permuteq}(\mathbf{rows},\mathbf{q})\land\texttt{diag\_clear}(\mathbf{q}).
\end{equation}
该写法更符合逻辑编程范式，但统一与回溯开销较大。

\section{M3：纯逻辑求解（自定义不等关系）}

\subsection{关系设计}

考虑到 Kanren 0.2.3 未提供可直接调用的 \texttt{neq}，脚本中补充了两类关系：
\begin{itemize}
  \item \texttt{neq\_rel(left, right)}：两项不相等；
  \item \texttt{diagonal\_neq\_rel(left, right, distance)}：两项差值绝对值不等于列距。
\end{itemize}

这两类关系在“按列赋值”过程中即时触发，可比“先生成再过滤”更早剪掉无效状态。

\subsection{关键代码}

\begin{lstlisting}[language=Python,caption={自定义不等关系实现},label=lst:custom-neq-rel]
def neq_rel(left, right):
    def goal(subst):
        left_value = reify(left, subst)
        right_value = reify(right, subst)
        if isvar(left_value) or isvar(right_value):
            return (subst,)
        if left_value == right_value:
            return ()
        return (subst,)

    return goal


def diagonal_neq_rel(left, right, distance: int):
    def goal(subst):
        left_value = reify(left, subst)
        right_value = reify(right, subst)
        if isvar(left_value) or isvar(right_value):
            return (subst,)
        if abs(left_value - right_value) == distance:
            return ()
        return (subst,)

    return goal
\end{lstlisting}

\section{M4 与 M5：DFS / BFS 图搜索}

二者都把“部分棋局”视为节点，“在下一列放置合法皇后”视为扩展操作。差异只在于前沿容器：
\begin{itemize}
  \item DFS 使用栈，优先纵向深入；
  \item BFS 使用队列，按层逐步展开。
\end{itemize}

\begin{algorithm}
  \caption{统一搜索框架（DFS/BFS）}\label{algo:search-framework}
  \begin{algorithmic}[1]
    \Require 棋盘规模 $N$，容器类型 $C$（栈或队列）
    \Ensure 全部合法解集合 $S$
    \State 将空状态 $()$ 放入 $C$
    \State $S \gets \varnothing$
    \While{$C$ 非空}
    \State 取出状态 $state$
    \If{$|state| = N$}
    \State 将 $state$ 加入 $S$
    \Else
    \For{每个候选行 $r \in [0, N-1]$}
    \If{$\text{safe}(state, r)$}
    \State 将 $state \oplus r$ 放入 $C$
    \EndIf
    \EndFor
    \EndIf
    \EndWhile
    \State \Return $S$
  \end{algorithmic}
\end{algorithm}

\section{M6：位运算回溯}

M6 使用位掩码维护列冲突与两条对角线冲突。对当前行而言，可放位置集合为
\begin{equation}
  \texttt{available}=\texttt{full\_mask}~\&~\sim(\texttt{col\_mask}|\texttt{main\_diag\_mask}|\texttt{anti\_diag\_mask}).
\end{equation}
每次提取最低位可行位置并递归更新掩码，即可完成搜索。

完整实现已放在附录的完整脚本清单中（见附录代码列表），正文不再重复展开整段函数。

\section{解集一致性验证}

六种方法都得到 92 个合法解；最小字典序解一致为
\begin{equation}
  (0,4,7,5,2,6,1,3),
\end{equation}
最大字典序解一致为
\begin{equation}
  (7,3,0,2,5,1,6,4).
\end{equation}
因此可确认不同实现在语义上等价，满足横向性能比较前提。
